\section{CEP Attribute Analysis and Resolution of Conflicts}\label{sec:cep}
\subsection{Attributes P1--P7}
\textbf{P1, depth of knowledge.} The solution combines distributed storage and processing (partitioned Parquet, broadcast joins, adaptive skew handling), data integration and quality engineering, surveillance statistics (EARS, Farrington-style seasonal baselines, CUSUM), count-time-series regression with overdispersion, proper scoring rules for interval forecasts, linear optimisation with equity constraints, visual encoding and interaction design, and statistical disclosure control.
\textbf{P2, conflicting requirements.} Five explicit conflicts (C1--C5) are resolved below with experiments rather than opinion.
\textbf{P3, depth of analysis.} Detection thresholds, forecast lags, the equity floor and the aggregation level are not given by any standard; each was chosen from a threshold sweep, an ablation or a scenario grid and compared with a baseline (Section~\ref{sec:eval}).
\textbf{P4, unfamiliar issues.} Reporting delays up to 21 days (and longer after back-filling), a facility that goes silent and back-fills, duplicate submissions, cross-facility duplicate patients, invalid values, skew towards Dhaka, and changes of reporting practice. The pipeline handles each explicitly: a checksum ledger and partition merge for late data, an artefact guard in detection, quarantine with reason codes, and a patient de-duplication window.
\textbf{P5, beyond standard codes.} No standard converts forecasts into an allocation plan. We combined WHO IDSR~\cite{idsr} surveillance principles, ICD-10 coding for disease groups and data-minimisation principles with our own conversion assumptions (kept in one versioned file) and a transparent optimisation model.
\textbf{P6, stakeholders.} Section~\ref{sec:cep} lists seven stakeholder groups; role-based views, uncertainty bands and an explanation of alpha respond to their specific concerns.
\textbf{P7, interdependence.} Quality affects detection (report completeness feeds the artefact guard), forecasts and uncertainty feed demand and the optimiser, and all of them feed the dashboard. A change of the epi-week definition or a late record propagates through all stages, which is why every stage is idempotent and tested end to end.

\subsection{Conflicts}
\textbf{C1, sensitivity versus false alarms.} We swept thresholds for three detectors and plotted sensitivity against false alert episodes per district-year (Fig.~\ref{fig:det}). The operating point maximises sensitivity subject to at most one false alert episode per district-year, a limit chosen because a district officer receiving more than about one spurious alert per year would stop reading them. @@C1_RESULT@@
\textbf{C2, granularity versus privacy.} The dashboard works at district level, never at facility or patient level. Counts of one to four are suppressed everywhere, including model estimates, and suppression is complementary: if a single cell of a group is hidden while the group total is shown, the next smallest cell is hidden too, so a hidden cell cannot be recovered by subtraction. Property tests over 5000 random groups confirm that no displayed value lies between one and four and no single hidden cell is recoverable. District officers see only their own district in detail.
\textbf{C3, efficiency versus equity.} The optimiser minimises cost-weighted unmet demand subject to an equity floor $\alpha$ (every district receives at least a fraction $\alpha$ of its population share of each resource). Over five budget levels and five values of $\alpha$ we plot unmet demand against the coverage of the worst-served district (Fig.~\ref{fig:eq}). @@C3_RESULT@@
\textbf{C4, freshness versus cost.} Decisions (allocation, campaigns) are weekly, reporting delays are up to three weeks, and counts for the latest weeks are incomplete, so a daily micro-batch of the analytics is sufficient. For the first signal of the week we added a Structured Streaming path that raises a provisional alert when the count already received for the current report week exceeds the full-week threshold. @@C4_RESULT@@
\textbf{C5, accuracy versus interpretability.} We use a negative-binomial regression whose coefficients officials can read (recent level, annual harmonics, rainfall and temperature lagged four to eight weeks) and compare it with the seasonal-naive baseline and with the same model without weather. @@C5_RESULT@@
